%=====================================================================
\chapter{User Account Management}\label{ch:accounts}
%=====================================================================
\locator{4}{Part IV --- Money and Identity}

\begin{outcomes}
By the end of this chapter you will be able to:
\begin{tight}
  \item \textbf{Explain} why a password is stored as a deliberately slow hash,
        and \textbf{name} what the cost factor buys.
  \item \textbf{Measure} the cost factor against throughput, and
        \textbf{choose} one for a stated shop.
  \item \textbf{Demonstrate} user enumeration by timing, and \textbf{explain}
        why a stronger hash makes it worse.
  \item \textbf{Evaluate} a lockout policy, and \textbf{identify} what it
        enables as well as what it prevents.
  \item \textbf{Design} a password reset flow that does not become the weakest
        way in.
\end{tight}
\end{outcomes}

\begin{prereq}
\chref{basket}: sessions, cookies, and that the server decides what is true.
\chref{gateway}: measuring a distribution rather than a single value. Appendix~A
--- every exercise here runs against your own container.
\end{prereq}

\begin{keyterms}
authentication $\bullet$ password hash $\bullet$ bcrypt $\bullet$ cost factor
$\bullet$ salt $\bullet$ rainbow table $\bullet$ credential stuffing $\bullet$
user enumeration $\bullet$ timing side channel $\bullet$ lockout $\bullet$
rate limiting $\bullet$ password reset $\bullet$ account takeover
\end{keyterms}

\section{Why a Password Is Stored Slowly on Purpose}

Every other piece of engineering in this course tries to make something faster.
Password storage is the one place where the software is deliberately slowed
down, and understanding why is the whole of the chapter.

\begin{definitionbox}[Hash, Salt, Cost Factor]
A \term{password hash} is a one-way function of the password. The shop stores
the output and never the password, so that a stolen database does not hand the
attacker the passwords.

\smallskip
A \term{salt} is a random value mixed in, different for every user, so that two
customers with the same password get different hashes --- and so that an
attacker cannot precompute a table of hashes once and look everybody up.

\smallskip
A \term{cost factor} says how much work the function does. Bcrypt's is a power
of two: raising it by one \textbf{doubles} the time. That is a cost to your
server once per login and a cost to the attacker once per \emph{guess}, and
the attacker is making billions of them.
\end{definitionbox}

\begin{examplebox}[What this shop actually does]
\begin{tight}
  \item Passwords are stored with PHP's \texttt{password\_hash(\dots,
        PASSWORD\_DEFAULT)}, which here produces a \texttt{\$2y\$10\$} hash ---
        \textbf{bcrypt at cost 10}, PHP's default and unchanged since 2013.
  \item The column is \texttt{varchar(255)}, which is correct: it leaves room
        for whatever algorithm PHP's default becomes next.
  \item Version 3 of this platform used SHA-1 with a salt, which is far too
        fast to be a password hash at all. Version 4 moved to
        \texttt{password\_hash}. \textbf{When you inherit a shop, check which
        one it is before anything else.}
  \item The minimum password length is \textbf{6 characters}, and none of the
        complexity settings --- uppercase, lowercase, number, symbol --- is
        set at all, so none is enforced.
\end{tight}
\end{examplebox}

\begin{measured}[The cost factor against the clock]
From \texttt{code/m15\_accounts.py --cost}, timed by PHP inside the shop
container on the machine in Appendix~A. ``Logins per second'' is one core doing
nothing else.

\rowcolors{2}{white}{lightbg}
\begin{longtable}{@{}L{2.2cm}L{2.4cm}L{2.6cm}L{3.0cm}L{3.2cm}@{}}
\toprule
\rowcolor{primary!15}
\textbf{Cost} & \textbf{Hash} & \textbf{Verify} & \textbf{Logins/sec} &
\textbf{Work vs cost 8} \\
\midrule
\endhead
8 & 13.6\,ms & 13.7\,ms & 73 & 1$\times$ \\
9 & 29.6\,ms & 26.8\,ms & 37 & 2$\times$ \\
\textbf{10} & \textbf{52.8\,ms} & \textbf{60.3\,ms} & \textbf{17} &
4$\times$ \\
11 & 110.2\,ms & 103.0\,ms & 10 & 8$\times$ \\
12 & 217.5\,ms & 218.9\,ms & 5 & 16$\times$ \\
13 & 440.7\,ms & 448.2\,ms & 2 & 33$\times$ \\
14 & 846.8\,ms & 873.8\,ms & 1 & 64$\times$ \\
\bottomrule
\end{longtable}

\smallskip
\textbf{The doubling is exact}, which is what a well-behaved cost factor looks
like: 13.7, 26.8, 60.3, 103, 219, 448, 874. Nothing else in this course
produces a table that clean, because nothing else in this course is a cost that
somebody designed on purpose.

\smallskip
\textbf{The shop ships at cost 10: 60\,ms per login, 17 logins per second on
one core.} A shop with 17 simultaneous login attempts is a shop in trouble for
other reasons, so cost 10 is not a throughput problem --- but it is also
twelve-year-old advice, and hardware has not stood still.

\smallskip
\textbf{Cost 12 is the usual current recommendation}: 219\,ms, 5 logins per
second. For a shop taking a few hundred orders a day that is free, and it makes
every stolen hash sixteen times more expensive to attack than cost 8.
\textbf{Raising it is one line, and the next section is the reason to think
first.}
\end{measured}

\section{What the Login Form Tells an Attacker Who Has No Password}

An attacker with a list of email addresses and no passwords still wants
something from your login form: \emph{which of these addresses have accounts
here}. That list is worth money on its own, and it makes every later attack
cheaper.

\begin{measured}[A registered email answers more slowly, and measurably so]
From \texttt{code/m15\_accounts.py --enumerate}: fifteen login attempts with a
deliberately wrong password, against an address that exists and one that does
not, timed from outside the container as an attacker would.

\smallskip
\begin{tabular}{@{}L{4.6cm}L{3.0cm}L{2.6cm}L{3.0cm}@{}}
& \textbf{Median} & \textbf{Fastest} & \textbf{Slowest} \\
\textbf{Registered email} & \textbf{88.6\,ms} & 70.9\,ms & 142.7\,ms \\
\textbf{Unknown email} & \textbf{23.0\,ms} & 15.2\,ms & 94.8\,ms \\
\end{tabular}

\smallskip
\textbf{A gap of 65.5\,ms --- 74\% of the slower time --- visible from
outside, over the network, without a single correct password.}

\smallskip
\textbf{Where the 65\,ms comes from is the uncomfortable part.} It is almost
exactly the cost of one bcrypt verify at cost 10, measured in the previous
section as 60.3\,ms. A wrong password on a real account must be checked against
the stored hash; a wrong \emph{email} matches no row, so there is nothing to
check and the request returns early.

\smallskip
\textbf{So the leak is produced by the security measure.} Raise the cost factor
to 12 for the excellent reasons above, and the gap grows from 65\,ms to roughly
220\,ms --- four times louder, trivially measurable over any network. A change
that makes stolen hashes sixteen times harder to crack makes your customer list
four times easier to enumerate.

\smallskip
\textbf{The fix is one line and it is not obvious:} when the email matches no
customer, verify the supplied password against a \emph{dummy hash} anyway and
discard the result. The request then costs the same whether the account exists
or not. It feels like wasted work because it is wasted work --- that is the
point.
\end{measured}

\begin{conceptbox}[Timing is not the only way this leaks]
Closing the timing channel is necessary and not sufficient. The same
information escapes through:

\begin{tight}
  \item \textbf{Different messages.} ``No such account'' versus ``Wrong
        password'' states it outright. Use one message for both.
  \item \textbf{Registration.} ``That email is already registered'' is the same
        disclosure through a different form, and it is harder to remove,
        because the user genuinely needs to know. Send an email instead of
        answering on the page.
  \item \textbf{Password reset.} ``We have sent you a link'' must be shown
        whether or not the address exists.
  \item \textbf{The lockout itself.} See below --- a locked account answers
        differently from one that was never there.
\end{tight}

\smallskip
\textbf{Enumeration is rarely the attack. It is the reconnaissance that makes
the attack affordable}, and that is why it is worth closing even though no
single leak looks serious.
\end{conceptbox}

\section{Lockout: What It Prevents, and What It Enables}

\begin{measured}[Five wrong guesses, then what?]
From \texttt{code/m15\_accounts.py --lockout}. The shop's
\texttt{config\_login\_attempts} is \textbf{5}.

\rowcolors{2}{white}{lightbg}
\begin{longtable}{@{}L{3.0cm}L{2.6cm}L{3.2cm}L{5.0cm}@{}}
\toprule
\rowcolor{primary!15}
\textbf{Attempt} & \textbf{Time} & \textbf{Recorded} & \textbf{Result} \\
\midrule
\endhead
1--5 & 70--81\,ms & 1, 2, 3, 4, 5 & rejected, hash verified each time \\
6 & \textbf{16\,ms} & 5 & \textbf{locked} \\
7 & \textbf{15\,ms} & 5 & locked \\
\textbf{correct password} & \textbf{17\,ms} & 5 & \textbf{still refused} \\
\bottomrule
\end{longtable}

\smallskip
\textbf{It works.} After five failures the account stops answering, for an
hour, keyed on the pair (email, IP address).

\smallskip
\textbf{Two things the table shows that the setting does not.}

\smallskip
\textbf{First, the correct password is refused too.} That is right --- a lock
that the attacker's next guess could open would be no lock --- but it means
\textbf{anyone who knows a customer's email address can lock them out of their
own account at will}, with five wrong guesses, for an hour, repeatedly. Lockout
converts a password-guessing attack into a denial-of-service one. It is the
right trade for most shops and it is a trade, not a free win.

\smallskip
\textbf{Second, the locked response is four times faster}, because it never
reaches the hash. The lockout is itself an enumeration channel: 16\,ms says
``this account exists and is locked'', which is more than the 23\,ms of an
unknown address was giving away.

\smallskip
\textbf{And the key is (email, IP).} An attacker with a few hundred addresses
has effectively no limit, while a school or an office behind one NAT shares a
single counter --- so the mechanism is weakest against the attacker it is for
and strongest against the customers it is not.
\end{measured}

\begin{alertbox}[Where the real risk is, and it is not your password form]
Customers reuse passwords. An attacker who has a list of email and password
pairs stolen from somewhere else tries them on your shop --- \term{credential
stuffing} --- and some of them work, because the passwords are genuinely
correct. Nothing in this chapter prevents that: the hash is fine, the lockout
counts five failures and these are not failures.

\smallskip
What helps is \textbf{not} another password rule. It is: checking new passwords
against lists of known-breached ones; offering a second factor; rate limiting by
IP and by address across the whole site rather than per account; and noticing
the pattern --- many accounts, one address, high success rate --- which is a
\chref{dashboard} problem rather than a login problem.

\smallskip
\textbf{Meanwhile this shop's minimum password length is six characters with no
complexity requirement at all.} That is the setting to change first, and it is
in \adminpath{System, Settings, Edit, Option}.
\end{alertbox}

\begin{conceptbox}[The reset flow is the other front door]
An account is only as secure as the cheapest way to get into it, and that is
usually the password reset.

\begin{tightnum}
  \item \textbf{The token must be unguessable and single-use}, and must expire
        in minutes rather than days.
  \item \textbf{It must not be in the URL of a page that loads third-party
        scripts}, which leak it in the \texttt{Referer} header.
  \item \textbf{Resetting must invalidate existing sessions}, or an attacker
        who was already in stays in after the customer ``fixes'' it.
  \item \textbf{The response must not reveal whether the address exists}, which
        is the enumeration problem again, in its easiest-to-fix form.
  \item \textbf{The email itself is now the weakest link.} You have delegated
        your authentication to the customer's mail provider, and should know
        that you have done so.
\end{tightnum}
\end{conceptbox}

\begin{pitfall}
\begin{tight}
  \item \textbf{Choosing a cost factor without measuring it on the machine that
        will run it.} Cost 14 is one login per second.
  \item \textbf{Raising the cost factor without closing the timing channel},
        which makes enumeration four times easier while you congratulate
        yourself.
  \item \textbf{Different error messages for ``no account'' and ``wrong
        password''.}
  \item \textbf{Believing lockout is free.} It hands anyone a denial-of-service
        tool against any customer whose email they know.
  \item \textbf{Rate limiting only per account}, which does nothing against an
        attacker trying one password across ten thousand accounts.
  \item \textbf{Leaving the minimum password length at six with no complexity
        rules}, which is this shop's shipped state.
  \item \textbf{A reset token that does not expire}, or that survives being
        used.
  \item \textbf{Testing any of this against a shop that is not yours.} See
        Appendix~A.
\end{tight}
\end{pitfall}

\begin{lab}[Measure your own login, then make it quieter]
\begin{lstlisting}[language=bash]
# Everything here runs against YOUR container on localhost. Read the
# scope statement in Appendix A before you start.

# --- 1. Find out what your shop stores ------------------------------
cd code
docker compose exec db mariadb -ushop -pshoppw opencart \
  -e "SELECT email, LEFT(password,7) FROM oc_customer LIMIT 3;"
# The first seven characters name the algorithm and the cost.

# --- 2. Measure the cost factor -------------------------------------
python m15_accounts.py --cost

# --- 3. Choose one, and justify it in a sentence --------------------
# How many logins per second does your shop actually need? Work it
# out from orders per day, not from a feeling.

# --- 4. Predict, then measure the timing leak -----------------------
# Does a wrong password on a REAL account take longer than one on an
# address that does not exist? Predict yes or no, and by how much.
python m15_accounts.py --enumerate

# --- 5. Explain the number you got ----------------------------------
# Compare the gap with the verify time from step 2. They should
# match. Say why.

# --- 6. Watch the lockout, and find what it gives away --------------
python m15_accounts.py --lockout
# Two separate problems are visible in that table. Name both.

# --- 7. Tighten the policy ------------------------------------------
# Admin: System > Settings > Edit > Option. Raise the minimum
# password length and turn on the complexity requirements. Then
# register a new customer and confirm the rules are enforced.

# --- 8. Write the fix you cannot make in the admin panel ------------
# Three lines of pseudocode: how would you make a login to a
# non-existent account take the same time as one to a real account?

# --- 9. Clean up ----------------------------------------------------
python m15_accounts.py --clean
\end{lstlisting}

\smallskip
\textbf{What to notice.} Step 5. The timing leak is not a bug somebody left in
--- it is the direct, unavoidable consequence of doing the right thing with
passwords. Security measures interact, and sometimes they interact badly.
\end{lab}

\begin{checkpoint}
\begin{tightnum}
  \item Your shop runs bcrypt at cost 10 and takes 400 orders a day. Work out
        whether cost 12 is affordable, stating your assumptions about logins
        per order and peak concentration, and say what else you would have to
        change at the same time and why.
  \item A wrong password on a registered address took 88.6\,ms and on an
        unknown one 23.0\,ms. Explain the cause, state what an attacker does
        with that, and give the one-line fix together with the reason it feels
        wrong to write.
  \item A client asks you to remove the five-attempt lockout because customers
        complain about being locked out. Give the risk you are accepting, the
        risk you are removing, and a design that addresses both complaints
        better than either keeping or removing it.
\end{tightnum}
\end{checkpoint}

\begin{chaptersummary}
\begin{tight}
  \item A password hash is slow on purpose. The salt defeats precomputation;
        the cost factor sets how expensive each guess is, and each step doubles
        it for you and for the attacker.
  \item This shop uses \texttt{password\_hash} with bcrypt at cost 10, in a
        \texttt{varchar(255)} column. Version 3 used salted SHA-1; check which
        one any shop you inherit is using before anything else.
  \item Measured: cost 8 to 14 ran 13.7, 26.8, 60.3, 103, 219, 448 and
        874\,ms --- an exact doubling. Cost 10 gives 17 logins per second on
        one core; cost 12, the usual current recommendation, gives 5.
  \item Measured: a wrong password on a registered address took 88.6\,ms and on
        an unknown address 23.0\,ms --- a 65\,ms gap, visible over the network,
        and almost exactly one bcrypt verify.
  \item The leak is caused by the defence. Raising the cost factor to 12 would
        make the gap roughly 220\,ms. The fix is to verify against a dummy hash
        when no account matches, so that both paths cost the same.
  \item Enumeration also leaks through error messages, registration, password
        reset and the lockout response. It is reconnaissance, not the attack.
  \item Measured: five wrong guesses lock the account for an hour, keyed on
        (email, IP). The correct password is refused too --- so lockout hands
        an attacker a denial-of-service tool --- and the locked reply is four
        times faster, which is itself a disclosure.
  \item Credential stuffing defeats all of it, because the passwords are
        correct. The answers are breach-list checks, a second factor, site-wide
        rate limiting, and noticing the pattern.
  \item The password reset is the other front door: single-use, short-lived
        tokens; session invalidation; and the knowledge that you have delegated
        authentication to a mail provider.
\end{tight}
\end{chaptersummary}

\begin{reviewq}
  \item \textbf{(MCQ)} Raising bcrypt's cost factor from 10 to 12 multiplies
        the work by: (a)~1.2; (b)~2; (c)~4; (d)~16.
  \item \textbf{(MCQ)} The measured 65\,ms gap between a registered and an
        unknown email is closest to: (a)~network latency; (b)~one bcrypt
        verify; (c)~the database lookup; (d)~template rendering.
  \item \textbf{(Short)} Explain what a salt prevents, and why it does not need
        to be secret.
  \item \textbf{(Short)} Give two ways a shop leaks which email addresses are
        registered, other than by timing, and say how to close each.
  \item \textbf{(Short)} Explain why account lockout is a trade rather than an
        improvement, naming what it enables.
  \item \textbf{(Applied)} A shop's database is stolen. All passwords are
        bcrypt at cost 10. Explain what the attacker can and cannot do, what
        the shop must tell its customers, and what it should change --- and
        say why changing the cost factor now helps less than you might expect.
  \item \textbf{(Applied)} Design rate limiting for a login form that must stop
        credential stuffing without locking out legitimate customers. State
        what you count, over what window, keyed on what, what happens at the
        limit, and how a school computer lab with forty students behind one
        address is not broken by your design.
\end{reviewq}
